\section{Two-Dimensional Fourier Neural Operator for Thermoelastic Field Prediction}

\subsection{Model Formulation}

In the present work, the Fourier Neural Operator is considered as a two-dimensional field-to-field prediction model. The model receives a tensor of physical fields, optionally supplements it with spatial coordinate information, transforms the input into a higher-dimensional latent representation, processes this representation through a sequence of Fourier blocks, and then projects the result to the required physical output fields.

The computational domain is represented as a two-dimensional spatial grid:
\begin{equation}
D \subset \mathbb{R}^{2}.
\end{equation}
Each point of the domain is described by two coordinates:
\begin{equation}
\boldsymbol{x} = (x,y).
\end{equation}
Although the spatial grid is two-dimensional, the physical state at each grid point may contain several channels. In this work, the target field is represented as
\begin{equation}
s(\boldsymbol{x}, t)
=
\left[
T(\boldsymbol{x}, t),
u(\boldsymbol{x}, t),
v(\boldsymbol{x}, t)
\right],
\end{equation}
where \(T(\boldsymbol{x}, t)\) is the temperature field, and \(u(\boldsymbol{x}, t)\) and \(v(\boldsymbol{x}, t)\) are displacement-related field components extracted from the simulation data.

The model learns an operator
\begin{equation}
\mathcal{G}_{\theta}: \mathcal{A} \rightarrow \mathcal{U},
\end{equation}
where \(\mathcal{A}\) is the input field space, \(\mathcal{U}\) is the output field space, and \(\theta\) denotes the trainable parameters of the neural network. The learned mapping can be written as
\begin{equation}
\hat{s}
=
\mathcal{G}_{\theta}(a),
\end{equation}
where \(a\) is the input field and \(\hat{s}\) is the predicted thermoelastic field.

\subsection{Input and Output Tensor Format}

The input tensor follows the channel-first convention:
\begin{equation}
A
\in
\mathbb{R}^{B \times C_{in} \times N_x \times N_y},
\end{equation}
where \(B\) is the batch size, \(C_{in}\) is the number of input physical channels, and \(N_x\), \(N_y\) are the numbers of grid points in the two spatial directions.

The output tensor has the form
\begin{equation}
\hat{S}
\in
\mathbb{R}^{B \times C_{out} \times N_x \times N_y}.
\end{equation}
For the present thermoelastic prediction task, the number of output channels is
\begin{equation}
C_{out}=3,
\end{equation}
because the model predicts three physical fields:
\begin{equation}
T, \quad u, \quad v.
\end{equation}

If several previous time steps are used as input, they can be concatenated along the channel dimension. For example, if each time step contains \(C\) physical channels and \(m\) previous time steps are used, then
\begin{equation}
C_{in}=mC.
\end{equation}
In that case, the input tensor becomes
\begin{equation}
A
\in
\mathbb{R}^{B \times (mC) \times N_x \times N_y}.
\end{equation}
This format allows temporal information to be included without treating time as an additional Fourier dimension.

\subsection{Main Architecture of the FNO Model}

The two-dimensional Fourier Neural Operator is defined by several architectural choices: the number of retained Fourier modes, the number of Fourier blocks, the number of input and output physical channels, and the dimensions of the latent representations used inside the model. These choices determine how much spectral information is preserved, how deep the operator approximation is, and how much internal capacity the model has for representing thermoelastic dynamics.

For a two-dimensional FNO, the retained Fourier modes are represented as
\begin{equation}
K_1, \quad K_2,
\end{equation}
where \(K_1\) and \(K_2\) are the numbers of retained Fourier modes in the two spatial directions. Since the computational grid is two-dimensional, the spatial dimension of the operator is
\begin{equation}
d = 2.
\end{equation}

The model consists of three main stages:
\begin{equation}
\text{input}
\rightarrow
\text{lifting}
\rightarrow
\text{Fourier blocks}
\rightarrow
\text{projection}
\rightarrow
\text{output}.
\end{equation}

This can be written in operator form as
\begin{equation}
\mathcal{G}_{\theta}
=
Q
\circ
\mathcal{B}_{L-1}
\circ
\cdots
\circ
\mathcal{B}_{1}
\circ
\mathcal{B}_{0}
\circ
P,
\end{equation}
where \(P\) is the lifting transformation, \(\mathcal{B}_l\) is the \(l\)-th Fourier block, \(L\) is the number of Fourier blocks, and \(Q\) is the projection transformation.

\subsection{Coordinate Grid Encoding}

The model may use coordinate encoding by appending normalized spatial coordinates to the input features before the lifting layer. For a two-dimensional grid, this means that the coordinate-augmented input at each grid point is
\begin{equation}
\tilde{a}(\boldsymbol{x})
=
\left[
a(\boldsymbol{x}),
x,
y
\right].
\end{equation}

If the original number of input channels is \(C_{in}\), then after adding the spatial coordinates the effective number of channels passed to the lifting layer becomes
\begin{equation}
C_{in} + 2.
\end{equation}
More generally, the number of added coordinate channels is equal to the spatial dimension \(d\). For the present two-dimensional case,
\begin{equation}
d=2.
\end{equation}

Coordinate encoding is useful because it gives the neural operator explicit information about the spatial position of every grid point. This is especially important when the same physical value may have different meaning depending on its location in the computational domain.

\subsection{Lifting Layer}

Before applying Fourier blocks, the model transforms the input channels into a higher-dimensional latent representation. The input is first rearranged from channel-first format
\begin{equation}
[B, C_{in}, N_x, N_y]
\end{equation}
to channel-last format
\begin{equation}
[B, N_x, N_y, C_{in}],
\end{equation}
so that a pointwise linear transformation can be applied at each grid point.

If coordinate encoding is used, the input channel dimension becomes
\begin{equation}
C_{in} + 2.
\end{equation}
The first lifting transformation is
\begin{equation}
v^{(1)}(\boldsymbol{x})
=
P_1
\left(
\tilde{a}(\boldsymbol{x})
\right),
\end{equation}
where
\begin{equation}
P_1:
\mathbb{R}^{C_{in}+2}
\rightarrow
\mathbb{R}^{C_{lift}}.
\end{equation}
Here, \(C_{lift}\) denotes the intermediate lifting dimension.

If an additional lifting stage is used, the representation is further mapped to the main hidden dimension:
\begin{equation}
v_0(\boldsymbol{x})
=
P_2
\left(
v^{(1)}(\boldsymbol{x})
\right),
\end{equation}
where
\begin{equation}
P_2:
\mathbb{R}^{C_{lift}}
\rightarrow
\mathbb{R}^{C_{mid}}.
\end{equation}
Here, \(C_{mid}\) denotes the main latent channel dimension of the Fourier blocks.

If the intermediate lifting stage is omitted, the input is mapped directly to \(C_{mid}\) using one pointwise linear transformation. Thus, the lifting stage prepares the latent field
\begin{equation}
v_0
\in
\mathbb{R}^{B \times C_{mid} \times N_x \times N_y},
\end{equation}
which is then processed by the sequence of Fourier blocks.

\subsection{Boundary Extension Before Fourier Blocks}

Before the sequence of Fourier blocks, the latent representation may be temporarily extended along the spatial boundaries:
\begin{equation}
v_0^{ext}
=
\operatorname{Ext}(v_0).
\end{equation}

This operation may be useful when spectral operations are applied to spatial data, because it can reduce boundary-related artifacts and make the internal spectral processing more stable for non-periodic physical fields. After all Fourier blocks are applied, the extended part is removed:
\begin{equation}
v_L
=
\operatorname{Crop}(v_L^{ext}).
\end{equation}

Thus, the boundary extension is not part of the physical prediction itself; it is an internal numerical step used during the forward pass of the neural network.

\subsection{Fourier Blocks}

The main processing part of the model is a sequence of Fourier blocks:
\begin{equation}
v_0
\rightarrow
v_1
\rightarrow
\cdots
\rightarrow
v_L.
\end{equation}
The number of such blocks is denoted by \(L\).

Each Fourier block contains a spectral convolution branch and a local convolution branch. The block can be represented as
\begin{equation}
v_{l+1}
=
\sigma
\left(
\mathcal{S}_l(v_l)
+
\mathcal{C}_l(v_l)
\right),
\end{equation}
where \(\mathcal{S}_l\) is the spectral convolution operation, \(\mathcal{C}_l\) is the local convolution operation, and \(\sigma\) is the activation function.

For a two-dimensional model, the local convolution branch can be expressed as a two-dimensional convolution:
\begin{equation}
\mathcal{C}_l
=
\operatorname{Conv}_{3 \times 3}.
\end{equation}
This local convolution preserves the spatial resolution because the convolution is applied with appropriate spatial padding.

The spectral branch is responsible for learning transformations in Fourier space, while the local convolution branch provides additional local spatial processing in the physical domain. Their outputs are added before applying the activation function.

\subsection{Spectral Convolution}

The spectral convolution layer is the central operation of the Fourier Neural Operator. Let the input to the spectral convolution be
\begin{equation}
v_l
\in
\mathbb{R}^{B \times C_{in}^{(l)} \times N_x \times N_y}.
\end{equation}
The model first applies a two-dimensional Fourier transform over the spatial dimensions:
\begin{equation}
\hat{v}_l
=
\mathcal{F}(v_l).
\end{equation}
For the current case, the transform is applied over the two spatial dimensions \(N_x\) and \(N_y\).

The Fourier representation is complex-valued:
\begin{equation}
\hat{v}_l
=
\hat{v}_{l,\mathrm{real}}
+
i\hat{v}_{l,\mathrm{imag}}.
\end{equation}
Therefore, the spectral operation uses complex-valued multiplication between Fourier coefficients and trainable spectral weights.

For the retained Fourier modes, the spectral transformation can be written as
\begin{equation}
\hat{z}_l(k_1,k_2)
=
R_l(k_1,k_2)
\hat{v}_l(k_1,k_2),
\end{equation}
where \(R_l(k_1,k_2)\) is a learnable complex-valued weight matrix for the frequency mode \((k_1,k_2)\).

The set of retained modes can be written as
\begin{equation}
\mathcal{M}_K
=
\left\{
(k_1,k_2)
:
0 \leq k_1 < K_1,
\quad
0 \leq k_2 < K_2
\right\}.
\end{equation}
Only the selected low-frequency modes are directly transformed by learned weights. All other modes remain zero in the output Fourier tensor:
\begin{equation}
\hat{z}_l(k_1,k_2)
=
0,
\qquad
(k_1,k_2) \notin \mathcal{M}_K.
\end{equation}

After spectral multiplication, the inverse Fourier transform is applied:
\begin{equation}
z_l
=
\mathcal{F}^{-1}
\left(
\hat{z}_l
\right).
\end{equation}
The output is converted back to the real spatial domain:
\begin{equation}
z_l
\in
\mathbb{R}^{B \times C_{out}^{(l)} \times N_x \times N_y}.
\end{equation}

\subsection{Complex Multiplication in Fourier Space}

The spectral weights are complex-valued. If the input Fourier coefficient is
\begin{equation}
\hat{v}
=
a + ib,
\end{equation}
and the spectral weight is
\begin{equation}
R
=
c + id,
\end{equation}
then their product is
\begin{equation}
R\hat{v}
=
(c+id)(a+ib).
\end{equation}
Expanding this expression gives
\begin{equation}
R\hat{v}
=
(ac-bd)
+
i(ad+bc).
\end{equation}

Therefore, the real and imaginary parts of the output are
\begin{equation}
\operatorname{Re}(R\hat{v})
=
ac-bd,
\end{equation}
\begin{equation}
\operatorname{Im}(R\hat{v})
=
ad+bc.
\end{equation}
This is the mathematical operation performed inside the spectral convolution layer. In tensor form, this multiplication is applied across input channels and retained Fourier modes, producing output channels in the spectral domain.

\subsection{Factorized Spectral Weights}

The spectral convolution may store spectral weights either in a dense form or in a compressed factorized form. The dense version stores full real and imaginary weight tensors:
\begin{equation}
R_l
\in
\mathbb{C}^{C_{in}^{(l)} \times C_{out}^{(l)} \times K_1 \times K_2}.
\end{equation}

A factorized representation reduces the number of trainable parameters by representing a full spectral weight tensor through a smaller core tensor and factor matrices. Conceptually, this can be written as
\begin{equation}
R_l
\approx
\mathcal{T}
\times_1 A_1
\times_2 A_2
\cdots
\times_n A_n,
\end{equation}
where \(\mathcal{T}\) is the core tensor and \(A_i\) are factor matrices. This representation can make the spectral layer more memory-efficient while preserving the main structure of the learned Fourier-domain transformation.

For the thesis description, the important point is that the spectral weights are trainable and may be stored either densely or in a compressed factorized form.

\subsection{Local Convolution Branch}

In addition to the spectral convolution, each Fourier block contains a local convolution branch. For the two-dimensional case, this branch is written as
\begin{equation}
\mathcal{C}_l(v_l)
=
\operatorname{Conv}_{3 \times 3}(v_l).
\end{equation}

The purpose of this branch is to process local spatial information directly in the physical domain. While the spectral branch captures global frequency-domain interactions, the local convolution branch captures neighboring spatial patterns. The output of the Fourier block is obtained by adding the spectral and local branches:
\begin{equation}
v_{l+1}
=
\sigma
\left(
\mathcal{S}_l(v_l)
+
\mathcal{C}_l(v_l)
\right).
\end{equation}

Here, \(\sigma\) is the activation function. In many Fourier Neural Operator models, the Gaussian Error Linear Unit is used:
\begin{equation}
\sigma = \mathrm{GELU}.
\end{equation}

\subsection{Pointwise Channel-Mixing Layer}

Some Fourier block variants may also include a pointwise channel-mixing layer. In the two-dimensional case, this operation can be represented using \(1 \times 1\) convolutions:
\begin{equation}
\mathcal{M}(v)
=
\operatorname{Conv}_{1 \times 1}^{(2)}
\left(
\sigma
\left(
\operatorname{Conv}_{1 \times 1}^{(1)}(v)
\right)
\right).
\end{equation}

The \(1 \times 1\) convolution does not mix neighboring spatial positions. Instead, it mixes information across channels at each grid point. This is equivalent to applying a small feed-forward neural network independently at every spatial location.

If this branch is included, the Fourier block can be written as
\begin{equation}
v_{l+1}
=
\sigma
\left(
\mathcal{S}_l(v_l)
+
\mathcal{C}_l(v_l)
+
\mathcal{M}_l(v_l)
\right),
\end{equation}
where \(\mathcal{M}_l\) denotes the pointwise channel-mixing branch. If it is not included, the practical structure of the block is the combination of spectral convolution, local convolution, and activation.

\subsection{Projection Layer}

After the sequence of Fourier blocks, the model projects the hidden representation to the output physical fields. Before projection, the tensor is rearranged from channel-first format
\begin{equation}
[B, C_{mid}, N_x, N_y]
\end{equation}
back to channel-last format
\begin{equation}
[B, N_x, N_y, C_{mid}].
\end{equation}

The first projection layer maps the hidden channels to an intermediate projection dimension:
\begin{equation}
q_1:
\mathbb{R}^{C_{mid}}
\rightarrow
\mathbb{R}^{C_{proj}},
\end{equation}
where \(C_{proj}\) denotes the intermediate projection dimension. After applying the activation function, the final linear layer maps the representation to the required number of output physical channels:
\begin{equation}
q_2:
\mathbb{R}^{C_{proj}}
\rightarrow
\mathbb{R}^{C_{out}}.
\end{equation}

Thus, the predicted field at each spatial point is
\begin{equation}
\hat{s}(\boldsymbol{x}, t)
=
q_2
\left(
\sigma
\left(
q_1(v_L(\boldsymbol{x}))
\right)
\right).
\end{equation}
For the present task,
\begin{equation}
C_{out}=3,
\end{equation}
and therefore
\begin{equation}
\hat{s}(\boldsymbol{x}, t)
=
\left[
\hat{T}(\boldsymbol{x}, t),
\hat{u}(\boldsymbol{x}, t),
\hat{v}(\boldsymbol{x}, t)
\right].
\end{equation}

Finally, the output is returned to channel-first format:
\begin{equation}
\hat{S}
\in
\mathbb{R}^{B \times C_{out} \times N_x \times N_y}.
\end{equation}
For the present three-channel thermoelastic task, this becomes
\begin{equation}
\hat{S}
\in
\mathbb{R}^{B \times 3 \times N_x \times N_y}.
\end{equation}

\subsection{Temporal Prediction Task}

The model can be used for one-step temporal prediction. If the known state at time \(t_n\) is
\begin{equation}
s(\boldsymbol{x}, t_n),
\end{equation}
then the model predicts the next state:
\begin{equation}
\hat{s}(\boldsymbol{x}, t_{n+1})
=
\mathcal{G}_{\theta}
\left(
s(\boldsymbol{x}, t_n)
\right).
\end{equation}

If several previous time steps are used, the input can be written as
\begin{equation}
a(\boldsymbol{x}, t_n)
=
\left[
s(\boldsymbol{x}, t_{n-m+1}),
s(\boldsymbol{x}, t_{n-m+2}),
\ldots,
s(\boldsymbol{x}, t_n)
\right],
\end{equation}
where \(m\) is the number of input time steps. The model then predicts
\begin{equation}
\hat{s}(\boldsymbol{x}, t_{n+1})
=
\mathcal{G}_{\theta}
\left(
a(\boldsymbol{x}, t_n)
\right).
\end{equation}

In tensor form, the temporal information is included by increasing the channel dimension:
\begin{equation}
A
\in
\mathbb{R}^{B \times (mC) \times N_x \times N_y}.
\end{equation}

\subsection{Supervised Training Objective}

The model can be trained in a supervised manner by comparing the predicted thermoelastic fields with reference fields obtained from simulation data. The loss function is not part of the FNO architecture itself, but belongs to the external training procedure. Let the prediction for sample \(j\) be
\begin{equation}
\hat{s}_j
=
\mathcal{G}_{\theta}(a_j),
\end{equation}
and let the corresponding reference field be
\begin{equation}
s_j.
\end{equation}

The mean squared error loss is defined as
\begin{equation}
\label{eq:mse}
\mathcal{L}_{MSE}
=
\frac{1}{N_sNC}
\sum_{j=1}^{N_s}
\sum_{i=1}^{N}
\sum_{c=1}^{C}
\left(
\hat{s}_{j,i,c}
-
s_{j,i,c}
\right)^2,
\end{equation}
where \(N_s\) is the number of samples, \(N\) is the number of spatial grid points, and \(C\) is the number of output channels. In the present case,
\begin{equation}
C = 3.
\end{equation}

The relative \(L^2\) loss is defined as
\begin{equation}
\mathcal{L}_{rel}
=
\frac{1}{N_s}
\sum_{j=1}^{N_s}
\frac{
\left\|
\hat{s}_j
-
s_j
\right\|_2
}{
\left\|
s_j
\right\|_2
}.
\end{equation}

If both criteria are used during training, the total loss can be written as
\begin{equation}
\mathcal{L}
=
\lambda_{MSE}\mathcal{L}_{MSE}
+
\lambda_{rel}\mathcal{L}_{rel},
\end{equation}
where \(\lambda_{MSE}\) and \(\lambda_{rel}\) are weighting coefficients.

\subsection{Normalization of Physical Channels}

Because the predicted channels have different physical meanings and may have different numerical scales, normalization is applied before training. For each channel \(c\), standard normalization is written as
\begin{equation}
\tilde{s}_{i,c}
=
\frac{
s_{i,c}
-
\mu_c
}{
\sigma_c
},
\end{equation}
where \(\mu_c\) is the mean value of the channel and \(\sigma_c\) is its standard deviation. The inverse transformation is
\begin{equation}
s_{i,c}
=
\tilde{s}_{i,c}\sigma_c
+
\mu_c.
\end{equation}

The normalization statistics are computed on the training set and reused for validation, testing, and inference. This ensures that the model receives inputs with comparable numerical scales and that the predicted normalized values can be converted back to physical units.

\subsection{Evaluation Metrics}

The trained model is evaluated by comparing predicted and reference fields. The evaluation MSE has the same form as the training-time mean squared error defined in Eq.~\eqref{eq:mse}, only without averaging over the sample batch when a single reference field is used at evaluation time.

The mean absolute error is
\begin{equation}
MAE
=
\frac{1}{NC}
\sum_{i=1}^{N}
\sum_{c=1}^{C}
\left|
\hat{s}_{i,c}
-
s_{i,c}
\right|.
\end{equation}

The relative error is
\begin{equation}
E_{rel}
=
\frac{
\left\|
\hat{s}
-
s
\right\|_2
}{
\left\|
s
\right\|_2
}.
\end{equation}

For channel-wise analysis, the relative error of channel \(c\) can be written as
\begin{equation}
E_c
=
\frac{
\left\|
\hat{s}_{c}
-
s_{c}
\right\|_2
}{
\left\|
s_c
\right\|_2
}.
\end{equation}

This metric is useful because the model predicts several physical fields at once, and each field may have a different range and prediction difficulty.

\subsection{Architectural and Training Parameters}

The most important architectural and training parameters of the two-dimensional FNO are:
\begin{equation}
K_1, K_2,
\quad
L,
\quad
C_{mid},
\quad
C_{lift},
\quad
C_{proj},
\quad
C_{out},
\quad
B,
\quad
\eta,
\end{equation}
where \(K_1\) and \(K_2\) are the retained Fourier modes, \(L\) is the number of Fourier blocks, \(C_{mid}\) is the hidden channel dimension, \(C_{lift}\) is the lifting dimension, \(C_{proj}\) is the projection dimension, \(C_{out}\) is the number of output channels, \(B\) is the batch size, and \(\eta\) is the learning rate.

The number of retained Fourier modes controls how much spectral information is used by the model. The number of Fourier blocks controls the depth of the operator approximation. The hidden channel dimension controls the capacity of the latent representation. The number of output channels defines how many physical fields are predicted by the projection layer. The batch size and learning rate affect training stability and computational cost.

\subsection{Summary}

In summary, the implemented FNO is a two-dimensional field-to-field prediction model. The input tensor is represented in channel-first format:
\begin{equation}
[B, C_{in}, N_x, N_y],
\end{equation}
and the output tensor has the form
\begin{equation}
[B, C_{out}, N_x, N_y].
\end{equation}
For the present thermoelastic configuration, this becomes
\begin{equation}
[B, 3, N_x, N_y].
\end{equation}
The model may append coordinate grid information, lifts the input to a hidden representation, applies a sequence of Fourier blocks, and then projects the latent field to the required output channels:
\begin{equation}
T, \quad u, \quad v.
\end{equation}

Each Fourier block in the main FNO path combines a spectral convolution branch with a local convolution branch. The spectral branch applies the Fourier transform, multiplies selected modes by learned complex-valued weights, and returns the result to the spatial domain using the inverse Fourier transform. The local convolution branch processes neighboring spatial information directly on the grid. The final prediction can be trained against simulation-generated reference fields using supervised losses such as mean squared error and relative \(L^2\) error.